\documentclass{article}
\usepackage[letterpaper,top=2cm,bottom=2cm,left=3cm,right=3cm,marginparwidth=1.75cm]{geometry}

\usepackage{graphicx} % Required for inserting images
\usepackage{booktabs}
\usepackage{caption}
\usepackage{subcaption}
\usepackage{authblk}
\usepackage{amsmath}
\usepackage{amssymb}
\usepackage[dvipsnames]{xcolor}
\usepackage[english]{babel}
\usepackage{biblatex}
\bibliography{latex/bibliography}
\usepackage[colorlinks=true, allcolors=blue]{hyperref}

\usepackage{verbatim}


% Change font
\usepackage{helvet}
\renewcommand{\familydefault}{\sfdefault}

\usepackage{xcolor}
\newcommand{\tkDE}[1]{\textcolor{magenta}{#1}} 
%\newcommand{\tkDE}[1]{\textcolor{magenta}{\textbf{David}: #1}}  

\newcommand{\diff}{\mathrm{d}}  

\title{Ethnicity-based System Dynamics Model of Measles Vaccinations in Birmingham and Solihull}
\author{David Ellis}

\date{}

\begin{document}
\maketitle

\begin{abstract}
    Abstract here.
\end{abstract}

\section{Introduction}

\tkDE{NHS MMR vaccine schedule.}

\tkDE{New MMRV vaccine.}

Factors influencing parental vaccine hesitancy \cite{Williams02092014}. MMR hesitancy in Europe meta analysis \cite{Tabacchi02072016}.

Trends in vaccine hesitancy can be politically motivated, however, these trends cannot be described by a simple left/right divide \cite{Babinska2026}.

\section{Methods}
\subsection{Stakeholder Engagement}
Following approach from Michail et. al \cite{michail2025evaluation}.

\subsection{Modelling}
Following approach from \cite{van2020explaining} using Stella Architect.

\section{Results}

\section{Discussion}


\appendix
\section{Contribution Extraction from Workshop Transcript}

Due to the length of the workshop, copilot was used to extract vocal contributions to supplement the written materials transcribed in Appendix.\ref{app:workshop-materials}. The full prompt given to Copilot was as follows:

\begin{verbatim}
I will upload two MS Teams transcript files (Part 1 and Part 2).

I need you to produce a full chronological extraction of contributions from both transcripts, 
ordered strictly by agenda items. The extractions should not include any names or otherwise 
identifiable information. Use the following agenda headings, in this order:

Introduction
10:00 Arrival
10:10 Welcome and project overview
10:20 Local epidemiology update (UKHSA)
10:30 Summary of key evidence on vaccination barriers
11:00 Brief introduction to system dynamics modelling
11:10 System dynamics interactive example (group task)
11:30 Break

Model Building
11:45 Base model introduction
11:55 Identify factors that influence key flows in the system (group task)
12:20 Discussion
12:30 Vote on most influential factors
12:45 Lunch Break

Interventions
13:30 Identify problematic flows in the system
13:40 Identify potential interventions to improve system behaviour (group task)
14:10 Discussion
14:30 Rank proposed interventions using an anonymous online vote
15:00 Finish

Instructions:

1. Go through both transcript files and extract *all* contributions in *chronological order*, 
rewriting them cleanly in full sentences.
2. Place each contribution under the correct agenda heading.
3. If the transcript text is unclear, truncated, or garbled, insert "[UNCLEAR]" at that 
point and summarise what is recoverable.
4. Do not omit any substantive content—include all contributions, group discussions, side 
comments, and facilitator guidance.
5. Produce the final output as a set of .txt files:
   - Part 0 (header/explanation)
   - Part 1 (Introduction section)
   - Part 2 (Model Building section)
   - Part 3 (Interventions section)
6. After generating Parts 1–3, compile them into a single combined file called 
"chronological_extraction_FULL.txt".

Please ask me if you need confirmation about formatting, but otherwise proceed without 
repeatedly checking after each step. Proceed without citations for transcript processing.
\end{verbatim}

\end{document}
